\section{The projective surface and its rational fiber}
\label{sec:surface}

We now turn the quotient of \Cref{prop:infinite-quotient} into a
fundamental-group image on a smooth projective surface.  The construction
must retain the distinguished fiber: its rational components will lift
compactly, while the fundamental group of their union will have infinite
image.

Recall the inputs.  The complement $U=Y\setminus D$ lies in a smooth
connected projective surface $Y$, the boundary $D$ has simple normal
crossings, and
\[
  \rho:\pi_1(U)\longrightarrow Q
\]
has dense image in an infinite pro-$2$ group.  Its restriction to a
punctured smooth fiber also has dense image.  Every component meridian of
$D$ has finite image, and the meridians of the distinguished fiber
$D_*$ have trivial image.  This fiber consists of an unmarked main
sphere and $24$ tails, each meeting the main sphere once and carrying
two horizontal markings.

\subsection{A finite cover that kills all compactification meridians}

At a crossing, two commuting meridians may have relations that neither
cyclic subgroup detects separately.  We therefore choose a finite
quotient that detects the entire local group at every boundary point.

\begin{lemma}[Finite local detection and descent]
\label{lem:local-detection}
Let $Y$ be a smooth connected projective complex surface, let $D$ be a
simple-normal-crossings divisor, and put $U=Y\setminus D$.  Suppose
$\rho:\pi_1(U)\to Q$ has dense image in a profinite group and sends
every component meridian of $D$ to an element of finite order.
There is a finite quotient $Q\twoheadrightarrow P_f$ with the following
properties.
\begin{enumerate}
  \item The quotient is injective on the image under $\rho$ of every
  sufficiently small local boundary-complement group.
  \item The connected finite cover
  $\pi:U'\to U$ defined by
  $\ker(\pi_1(U)\to P_f)$ extends to a projective morphism
  $p:X\to Y$, where $X$ is smooth and connected and
  $D_X=X\setminus U'$ is a simple-normal-crossings divisor.
  \item Writing $i:U'\hookrightarrow X$ for inclusion, there is a
  homomorphism $\rho_X:\pi_1(X)\to Q$ such that
  \begin{equation}
    \rho_X\circ i_* = \rho\circ\pi_*.
    \label{eq:rho-descent}
  \end{equation}
\end{enumerate}
Moreover, the resolution used to construct $X$ can be chosen to preserve
every open subset of the normalization on which the surface is smooth
and its reduced boundary has simple normal crossings.
\end{lemma}

\begin{proof}
We first find a finite quotient detecting all local boundary images,
then compactify and resolve the corresponding cover, and finally
check that the original quotient map kills every new boundary
meridian.

Choose a small coordinate polydisk $V$ about a point of $D$.  Its
boundary complement is
\[
  V\setminus D\simeq(\Delta^*)^r\times\Delta^{2-r},
  \qquad r\in\{1,2\}.
\]
After choosing an access path in $U$, let
$A_V\subset Q$ be the image of its fundamental group.  This group is
generated by $r$ commuting finite-order elements, so it is finite.
Only finitely many such images occur up to conjugacy: there is one
meridian type on the smooth part of each boundary component and one
local group at each crossing.  The smooth part of an irreducible
component, with its finitely many crossings removed, is connected;
moving the access path along it conjugates the local group.

For every nonidentity element in representatives of these finite
groups, choose an open normal subgroup of $Q$ that does not contain
it.  Their finite intersection $N$ is open and normal.  It satisfies
\begin{equation}
  N\cap gA_Vg^{-1}=\{1\}
  \qquad(g\in Q)
  \label{eq:local-detection}
\end{equation}
for every local group.  Set $P_f=Q/N$.  Density makes the map
$\pi_1(U)\to P_f$ surjective, so its kernel defines a connected
finite Galois cover $\pi:U'\to U$.

The cover has the complex structure obtained by lifting local charts.
It is algebraic and finite \emph{\'etale} by the Riemann existence
theorem for complex varieties
\citep[Expos\'e XII, Th\'eor\`eme 5.1]{SGA1}.  This theorem applies
to the nonproper surface $U$.  Normalize $Y$ in the function field of
$U'$ and denote the resulting morphism by $n:\overline Y\to Y$.
It is finite: complex varieties are Nagata, so finiteness of relative
normalization applies to the finite-type morphism $U'\to Y$
\citep[Tag 03GR]{StacksProject}.
Its restriction over $U$ is $U'$, because the latter is already normal.
Finite morphisms are projective, so $\overline Y$ is projective
\citep[Tag 0B3I]{StacksProject}.

Resolve $\overline Y$ by projective blowups, leaving its smooth locus
unchanged; strong resolution in characteristic zero has this
prescribed-open property \citep[pp.~822--823]{EncinasHauser2002}.
Next resolve the reduced boundary on the resulting smooth surface.
For curves on a smooth surface this can be done by point blowups:
resolve the singular component germs, decrease tangency intersection
multiplicities, and separate crossings of three or more branches
\citep[Tag 0BIC and its proof]{StacksProject}.
Choose these centers only where the boundary is not already simple
normal crossings.  Thus this additional operation preserves every
open subset on which the pair was already smooth with simple normal
crossings.  We obtain
\[
  X\xrightarrow{\ r\ }\overline Y
    \xrightarrow{\ n\ }Y,
  \qquad p=n\circ r,
\]
with $X$ smooth and projective and with $D_X$ a
simple-normal-crossings divisor.  The dense open set $U'$ is unchanged,
so $X$ is connected.

We prove the factorization, including the meridians of resolution
exceptional curves.  The inclusion $i$ is surjective on fundamental
groups: a loop can be perturbed off a real-codimension-two divisor.
Its kernel is normally generated by the component meridians of $D_X$.
Indeed, perturb a null-homotopy disk relative to its boundary so that
it avoids the finitely many crossings of $D_X$ and meets its smooth
part transversely.  Deleting small disks around the finitely many
intersection points expresses its boundary loop as a product of
conjugates of meridians and their inverses.

Let $E$ be any component of $D_X$, choose a smooth point $x\in E$
away from the other components, and take a small transverse disk
$T$ through $x$ with $T\setminus\{x\}\subset U'$.  Choose a boundary
polydisk $V$ about $p(x)$ downstairs.  By shrinking $T$, we may assume
$p(T)\subset V$.  Consequently the projection of a small meridian
in $T\setminus\{x\}$ lies in $V\setminus D$.  After including its
access path, its $Q$-image belongs to a conjugate of $A_V$.
This argument applies also when $E$ is exceptional and is contracted
to a boundary point.

The projected meridian has trivial image in $P_f$, since its lift is
a closed loop in $U'$.  Its $Q$-image therefore lies in both $N$ and
a conjugate of $A_V$, and is trivial by
\eqref{eq:local-detection}.  All generators of $\ker i_*$ are thus
killed by $\rho\circ\pi_*$.  This proves
\eqref{eq:rho-descent}.
\end{proof}

Apply \Cref{lem:local-detection} to the map supplied by
\Cref{prop:infinite-quotient}.  Retain its notation
$P_f,N,U',\overline Y,X,p,\rho_X$.  We next verify that the
normalization is already smooth along the distinguished fiber, so
that the resolution can preserve it.

\subsection{The distinguished fiber remains a union of spheres}

\begin{lemma}
\label{lem:rational-special-fiber}
Let $h:X\to C'$ be the composite of $p$ and $Y\to C'$.  The
resolution can be chosen so that the reduced fiber
\[
  Z=h^{-1}(c_*)_{\mathrm{red}}
\]
has only smooth rational irreducible components and ordinary
transverse crossings.
\end{lemma}

\begin{proof}
Every vertical meridian along $D_*$ is trivial in $Q$, and hence in
$P_f$, by \Cref{lem:distinguished-meridians}.  We describe the
normalization locally at every kind of point of $D_*$.  The
restriction of the Galois cover to a boundary-complement polydisk
is classified by its local meridian homomorphism to $P_f$; its
connected components correspond to cosets of the local image.

At a smooth unmarked point of $D_*$, the sole meridian is trivial.
Every connected local cover therefore extends as a copy of the
polydisk.  At a node between the main component and a tail, both
vertical meridians are trivial, and the same assertion holds.

At a horizontal--vertical intersection choose coordinates $(u,v)$
with the vertical curve $u=0$ and the horizontal marking $v=0$.
The $u$-meridian is trivial.  If the image of the $v$-meridian has
order $e$, the kernel of the local map $\Z^2\to P_f$ is
$\Z\times e\Z$.  Each connected local cover thus has the finite
normal extension
\begin{equation}
  (u,w)\longmapsto (u,v)=(u,w^e).
  \label{eq:distinguished-local-cover}
\end{equation}
This extension is smooth, and its reduced boundary is $uw=0$.
Analytification preserves finiteness and normality
\citep[Expos\'e XII, Propositions~2.1(vi) and 3.2(vi)]{SGA1}.
Thus the local extensions just described agree with $\overline Y$ by
uniqueness of finite normal extension
\citep[Expos\'e XII, Proposition~5.3]{SGA1}. The entire preimage of $D_*$
lies in the smooth simple-normal-crossings locus of the normalized
pair.  The resolution in \Cref{lem:local-detection} can be the
identity on this open neighborhood.



Consider now an irreducible component above the main sphere.
Its map to that sphere is unramified, including at the attaching
nodes, and there are no horizontal markings on the main component.
It is therefore a connected unramified cover of $\PP^1$, hence a
sphere mapping with degree one.  A component above a tail can branch
only at its two horizontal markings, and is unramified at the
attaching node.  Removing those two points downstairs and all their preimages upstairs
gives a connected finite cover of $\C^*$. Its covering subgroup is
$d\Z\subset\Z$.  Thus it is the covering $z\mapsto z^d$ of $\C^*$,
whose smooth compactification is $\PP^1$.

There are no additional components from resolution over $D_*$.
The local models show that the components just identified are
smooth and meet only in ordinary transverse crossings.  This
proves the assertion about $Z$.
\end{proof}

\subsection{Infinite image on the compact fiber}

Although $D_*$ is a tree of spheres, a degree-$d$ component above
one of its tails meets $d$ distinct degree-one lifts of the main
sphere: the attaching node has $d$ distinct preimages. The incidence
graph upstairs therefore need not be a tree. For a connected nodal
union of simply connected components, van Kampen identifies its
fundamental group with that of its dual graph, which has one vertex
for each component and one edge for each node. We now prove that
some connected component of $Z$ has infinite fundamental-group image
in $X$.

\begin{lemma}
\label{lem:infinite-special-image}
Let $Z=h^{-1}(c_*)_{\mathrm{red}}$ be the rational nodal fiber of
\Cref{lem:rational-special-fiber}. Some connected component $Z_0$ of $Z$ satisfies
\[
  \bigl|\im\bigl(\pi_1(Z_0)\longrightarrow\pi_1(X)\bigr)\bigr|
  =\infty.
\]
\end{lemma}

\begin{proof}
Let $F_s\subset U$ be a punctured smooth fiber, with $s\in B'$,
and choose access paths when comparing its group with $\pi_1(U)$.
Its image under $\rho$ is dense in $Q$ by
\Cref{prop:infinite-quotient}.  In particular its image in $P_f$
is all of $P_f$, so its inverse image $F'_s\subset U'$ is
connected.  Under the covering inclusion, its group is the
preimage of $N$ in $\pi_1(F_s)$.  Thus the image of
$\pi_1(F'_s)$ in $Q$ is
\[
  \rho\bigl(\pi_1(F_s)\bigr)\cap N.
\]
It is dense in $N$, because $N$ is open and the fiber image is
dense in $Q$.  The group $N$ is infinite: a finite open subgroup
would make its finite-index overgroup $Q$ finite.  Consequently
the image of $\pi_1(F'_s)$ in $Q$ is infinite.  By
\eqref{eq:rho-descent}, its image in $\pi_1(X)$ is infinite as
well.

We transfer this conclusion to $Z$ using an explicit
deformation-retract neighborhood.  Regard $X$ as a compact real analytic
manifold and $Z$ as a closed semianalytic subset.  This pair admits a
compatible finite triangulation
\citep[Section 3, Theorem 2, p.~464]{Lojasiewicz1964}.
Write it as a finite simplicial pair $(\mathcal K,\mathcal A)$,
and barycentrically subdivide so that $\mathcal A$ is a full
subcomplex: any simplex whose vertices belong to $\mathcal A$
belongs to $\mathcal A$.  For a point of $|\mathcal K|$, let $s_{\mathcal A}$
be the sum of its barycentric coordinates at vertices of
$\mathcal A$.  On the open neighborhood
\[
  V=\{s_{\mathcal A}>0\}
\]
retain those coordinates, divide them by $s_{\mathcal A}$, and
set all other coordinates to zero.  Fullness ensures that the
result belongs to $|\mathcal A|$.  The straight-line homotopy
inside each simplex gives a strong deformation retraction
$V\to Z$.  In particular the inclusion of $V$ into $X$ is
homotopic to a map through $Z$.

Properness of $h$ supplies the needed containment of nearby
fibers.  The set $h(X\setminus V)$ is closed in $C'$ and does
not contain $c_*$.  Hence every whole fiber over a sufficiently
small neighborhood of $c_*$ lies in $V$.  Choose $s\in B'$ in
this neighborhood.  Its connected punctured lifted fiber
$F'_s$ lies in one connected component of $V$; the deformation
retraction takes that component into a connected component
$Z_0$ of $Z$.  Therefore its inclusion-induced homomorphism to
$\pi_1(X)$ factors, up to a change-of-basepoint conjugation,
through $\pi_1(Z_0)$.  Since the former has infinite image, so
does the latter.
\end{proof}

\begin{proof}[Proof of \Cref{thm:main}]
The construction above gives a smooth connected projective surface
$X$.  By \Cref{lem:rational-special-fiber,lem:infinite-special-image},
it contains a connected nodal curve $Z_0$ with smooth rational
irreducible components and infinite image in $\pi_1(X)$.

\Cref{lem:rational-hull} now applies to $(X,Z_0)$.  In the simply
connected universal cover $\widetilde X$, every connected component
$W$ above $Z_0$ is a closed noncompact locally finite union of
compact rational curves.  For every $w\in W$, the holomorphic
hull of $\{w\}$ contains $W$ and is noncompact.  This proves all
the assertions of \Cref{thm:main}.
\end{proof}
